\section{The global Hall and direct-sum comparison}
\label{sec:elliptic-hall-comparison}

The tensor symmetry is decisive for the two geometric operations.
On the full spherical carrier, the ordinary symmetry is incompatible
with the Hall relation already in charge two. The correction is a
polynomial operator on the actual common-base model. Its invariant
tensors are the ordinary invariant tensors in
\eqref{gh:curve-cohomology}, which allows the full flag trace to be
identified with its symmetrizer.

\subsection{The diagonal in the Hall product}

For ordered blocks $A,B$ of labels define
\[
 K_{A,B}=\prod_{i\in A,\ j\in B}
              \left(1+\frac{\delta_{ij}}{z_j-z_i}\right).
\]
Here $\delta_{ij}$ is the class \eqref{gh:diagonal} on those two
curve factors. Let $\operatorname{Sh}(a,b)$ be the permutations of
$\{1,\ldots,a+b\}$ whose images of $1,\ldots,a$ are increasing and
whose images of $a+1,\ldots,a+b$ are increasing.
Permutations act by moving old factor $i$ to its new position.
Write $\tau_w$ for this simultaneous action, with its Koszul sign.
For $F\in\mathsf B_a\otimes\mathsf B_b$, the geometric Hall product is
\begin{equation}\label{gh:kernel-product}
 m_{a,b}(F)=
 \sum_{w\in\operatorname{Sh}(a,b)}\tau_w(FK_{A,B}),
 \qquad A=\{1,\ldots,a\},\quad B=\{a+1,\ldots,a+b\}.
\end{equation}
The right side is written in a rational localization, but its sum
is polynomial and equals the actual proper-trace product
\cite[\S3.3]{JindalLim2026}.
This polynomiality is part of the geometric comparison. With
permuted cohomological coefficients, it cannot be deduced merely by
calling a numerator anti-invariant.

For example, in charge $(1,1)$,
\begin{equation}\label{gh:binary-products}
 z*1=z_1+z_2-\delta_{12},\qquad
 1*z=z_1+z_2+\delta_{12},\qquad
 [z,1]_*=-2\delta_{12}\ne0.
\end{equation}
The elements $1,z$ here both have charge one; the algebra unit has
charge zero. In particular they are not commuting current
generators.

\begin{proposition}[The ordinary-symmetry obstruction]
\label{gh:ordinary-obstruction}
The actual graded Hall algebra
$\mathcal H=\bigoplus_{n\geq0}\mathsf B_n$, with
$\mathsf B_0=\mathbb Q$, admits no charge-preserving bialgebra
coproduct for the usual Koszul symmetry and its given unit.
\end{proposition}
\begin{proof}
In any connected nonnegatively charge-graded bialgebra, every
charge-one element is primitive. Hence the component of
$\Delta m$ in output charges $(1,1)$ is
$\operatorname{id}+\tau$ on $N\otimes N$.
But \eqref{gh:kernel-product} gives the quadratic relation
\[
 r=z_1-z_2+\delta_{12},\qquad m_{1,1}(r)=0.
\]
The ordinary swap fixes $\delta_{12}$, so
\[
 (\operatorname{id}+\tau)r=2\delta_{12}\ne0.
\]
This contradicts coproduct compatibility.
\end{proof}

In particular the direct-sum coalgebra of
Proposition~\ref{gh:direct-sum} is not an ordinary bialgebra
coproduct. On $z\otimes1$ its defect is
\[
 d_{1,1}(z*1)-(z_1+z_2)=-\delta_{12}.
\]
The defect $-\delta_{12}$ of this map and the relation obstruction
$2\delta_{12}$ above are different assertions. The latter excludes
every coproduct with the stated product, charge, unit and ordinary
symmetry. It does not exclude a different tensor recipient.

\subsection{A polynomial crossing over the common base}

On two ordered factors let $s_E$ exchange the exterior generators
with their Koszul signs, and let $s_z$ exchange only $z_1,z_2$.
Thus $\tau=s_Es_z$. Define
\begin{equation}\label{gh:polynomial-R}
 \partial_z=\frac{1-s_z}{z_1-z_2},\qquad
 R=\tau-\delta_{12}\partial_z .
\end{equation}
The divided difference is a polynomial operator: it is applied
coefficientwise to the Chern polynomial. Its degree is $-2$,
while multiplication by $\delta_{12}$ has degree two. No Chern
difference has been inverted in the definition of $R$.

\begin{proposition}\label{gh:R-identities}
The operator $R$ is a degree-zero $\Lambda_E$-linear chain map.
On all tensor powers its adjacent operators satisfy
\[
 R_i^2=1,\qquad
 R_iR_{i+1}R_i=R_{i+1}R_iR_{i+1},\qquad
 R_iR_j=R_jR_i\quad(|i-j|>1).
\]
It therefore gives an actual relative crossing on
$\mathcal A_1\otimes^{\mathbb L}_{\mathbb Q_{B_E}}\mathcal A_1$.
\end{proposition}
\begin{proof}
Both diagonal contractions kill $\delta_{12}$; all permutations
and divided differences commute with the diagonal action.
Thus $R$ is $\Lambda_E$-linear and commutes with
$uI_a+vI_b$ in \eqref{gh:koszul-model}.

For the involution, the polynomial identities are
\[
 s_z^2=1,\qquad s_z\partial_z+\partial_zs_z=0,\qquad
 \partial_z^2=0.
\]
On exterior coefficients,
$s_E^2=1$ and $s_E\delta_{12}=\delta_{12}s_E=\delta_{12}$ as
operators. Indeed $f-s_Ef$ belongs to the ideal generated by
$a_1-a_2,b_1-b_2$, which is annihilated by their exterior product.
Expanding $R^2$ gives the identity.

Here is the three-factor calculation. Permuting indices permutes
the classes $\delta_{ij}$, and
\[
 \delta_{12}\delta_{13}
 =\delta_{12}\delta_{23}
 =\delta_{13}\delta_{23}=:\delta_{123},\qquad
 \delta_{ij}^2=0.
\]
For instance, substitute
$a_1-a_3=(a_1-a_2)+(a_2-a_3)$ and the analogous equality for $b$.
Multiplication by $\delta_{12}$ kills the repeated differences,
leaving the claimed product. Multiplication by $\delta_{123}$
annihilates every pairwise difference, so every coefficient
permutation acts identically after that multiplication.

Write $s_1,s_2$ for the adjacent Chern permutations and
$\partial_1,\partial_2$ for their divided differences.
The one-correction terms in the two braid expansions agree by
\[
 \begin{gathered}
 s_1s_2\partial_1=\partial_2s_1s_2,\qquad
 s_1\partial_2s_1=s_2\partial_1s_2,\\
 \partial_1s_2s_1=s_2s_1\partial_2 .
 \end{gathered}
\]
All two-correction coefficients equal $\delta_{123}$, and their
polynomial operators agree by
\[
 \begin{split}
 s_1\partial_2\partial_1+\partial_1s_2\partial_1
                         +\partial_1\partial_2s_1
 ={}&s_2\partial_1\partial_2+\partial_2s_1\partial_2\\
    &+\partial_2\partial_1s_2 .
 \end{split}
\]
These identities follow by substituting
$\partial_i=(1-s_i)/(z_i-z_{i+1})$ and using
\[
 \frac1{(z_1-z_2)(z_1-z_3)}
 +\frac1{(z_2-z_3)(z_1-z_3)}
 =\frac1{(z_1-z_2)(z_2-z_3)}.
\]
The zero-correction terms agree by the permutation braid relation;
the three-correction terms vanish because a diagonal factor repeats.
This proves the braid equation on the polynomial module.
Disjoint operators commute. This is the elliptic specialization of
the crossing calculation in \cite[\S3.5]{JindalLim2026}, now on
the common-base action model.
\end{proof}

\begin{lemma}[The two invariant subspaces]\label{gh:invariants}
For every $n$,
\[
 T_n^{R(S_n)}=T_n^{S_n}=\mathsf B_n .
\]
Consequently
$e_n=\frac1{n!}\sum_{w\in S_n}R_w$ is a
$\Lambda_E$-linear idempotent with image $\mathsf B_n$.
\end{lemma}
\begin{proof}
For one adjacent pair, $s_z\partial_z=\partial_z$, and therefore
$\tau(\delta_{12}\partial_z f)=\delta_{12}\partial_z f$.
If $Rf=f$, then
$f-\tau f=-\delta_{12}\partial_z f$ is at once
$\tau$-anti-invariant and $\tau$-invariant. Over $\mathbb Q$ it
vanishes, so $\tau f=f$.
Conversely, if $\tau f=f$, then $s_zf=s_Ef$, and hence
$\delta_{12}(1-s_z)f=\delta_{12}(1-s_E)f=0$.
The Chern polynomial $z_1-z_2$ is a nonzero divisor on the free
module with exterior coefficients, so
$\delta_{12}\partial_zf=0$ and $Rf=f$.
Intersect these equalities over adjacent pairs. Averaging the
resulting symmetric-group representation proves the idempotent
assertion.
\end{proof}

The equality of invariant subspaces does not assert that $R=\tau$.
For example $R(z_1)=z_2-\delta_{12}$.
Both crossings have square one; erasing the diagonal term is not
detected by that equation.

\begin{lemma}[A restricted test for relative maps]
\label{gh:free-source}
The zero-differential action module $N^{\otimes n}$ is a direct sum
of shifted free $\Lambda_E$-modules for every $n>0$.
A degree-zero derived map from it to a zero-differential action
module is therefore determined by its underlying cohomology map.
In particular the binary Hall relation $m_{1,1}R=m_{1,1}$ can be
checked on its actual fibre cohomology.
\end{lemma}
\begin{proof}
The element $ab$ generates $\bigwedge(a,b)$ freely under the two
contractions, and each $z^k$ gives another shifted copy.
For a Hopf algebra, the diagonal action on a free module tensored
with any module is free. The change of variables uses
$\lambda\otimes v\mapsto\sum\lambda_{(1)}
\otimes\lambda_{(2)}v$; the inverse uses the antipode on the second
factor, with the graded tensor signs. Iteration proves the first
claim.

A zero-differential direct sum of shifted free modules is
K-projective: its Hom complex into an acyclic module is a product of
shifted acyclic rational complexes, hence acyclic. Derived maps
out of it are represented by strict chain maps. If the target also
has zero differential, the Hom differential is zero, so distinct
degree-zero maps are distinguished by their underlying cohomology
maps. Theorem~\ref{gh:simultaneous} supplies the actual action models
for the binary source and target. Formula
\eqref{gh:kernel-product} gives $m_{1,1}=1+R$ after inclusion into
$T_2$, and $R^2=1$ gives the asserted relation.
\end{proof}

This criterion uses its free source and zero-differential target.
It is not a faithfulness statement for the fibre functor on all
derived objects over $BE$.

\subsection{The tensor category generated by the crossing}

We specify the morphisms on which the new symmetry is natural.
Start with a $\mathbb Q$-linear strict monoidal category whose
objects are integers $n\geq0$, realized by $T_n$ with $T_0=\mathbb Q$.
Its endomorphisms at $n$ are the image of
$\mathbb Q[S_n]\to\operatorname{End}_{\Lambda_E}(T_n)$ under $R$,
and its morphism spaces between distinct integers are zero.
Tensor is concatenation, with the disjoint-block embeddings of the
symmetric groups. Let $\mathcal T_R$ be its additive idempotent
completion: objects are finite direct sums of images of these
idempotents, and maps are the corresponding matrix maps.

The block permutation exchanging an $a$-block and a $b$-block gives
a symmetry
$\beta_{a,b}:T_a\otimes T_b\to T_b\otimes T_a$.
Conjugation by this permutation exchanges the two block group
algebras, proving naturality for every allowed morphism.
The permutation identities give the two hexagons and
$\beta_{b,a}\beta_{a,b}=1$.
Thus $\mathcal T_R$ is a symmetric monoidal category with a
faithful monoidal realization in rational action complexes.
The realized symmetry is the block $R$-operator, rather than the
ambient Koszul swap. Passage to the corresponding relative
derived objects is monoidal; it defines no new symmetry on all
of $D(B_E;\mathbb Q)$.

The object $(n,e_n)$ realizes $\mathsf B_n$.
We allow charge-graded objects in $\mathcal T_R$, with convolution
tensor
$(V\otimes W)_n=\bigoplus_{a+b=n}V_a\otimes W_b$.
Each such sum is finite. For $F\in\mathsf B_a\otimes\mathsf B_b$
define
\begin{equation}\label{gh:abstract-operations}
 m^R_{a,b}(F)=\sum_{w\in\operatorname{Sh}(a,b)}R_wF,\qquad
 d^R_{a,b}:\mathsf B_{a+b}\hookrightarrow
                           \mathsf B_a\otimes\mathsf B_b .
\end{equation}
These are morphisms in $\mathcal T_R$. The first is a group-algebra
sum with image in $e_{a+b}$; the second is represented by $e_{a+b}$,
since $(e_a\otimes e_b)e_{a+b}=e_{a+b}$.

\begin{theorem}[The polynomial Hopf model]\label{gh:formal-hopf}
The operations \eqref{gh:abstract-operations} make
$\mathsf B=\bigoplus_{n\geq0}\mathsf B_n$ a connected Hopf
algebra in the specified charge-graded tensor category.
The unit is $\mathsf B_0=\mathbb Q$, the counit is projection onto
charge zero, and the antipode is
$S|_{\mathsf B_n}=(-1)^n\operatorname{id}$.
Its strict primitive subobject is $\mathsf B_1=N$.
As an algebra it is the braided symmetric algebra
\begin{equation}\label{gh:braided-symmetric}
 \operatorname{Sym}^R(N)=
 T(N)\big/\big\langle f-Rf:\ f\in N\otimes N\big\rangle ,
\end{equation}
where a tensor of length $n$ maps to
$\sum_{w\in S_n}R_w$ of that tensor.
\end{theorem}
\begin{proof}
Associativity of the shuffle product follows by identifying both
iterated sums with the interleavings of three ordered lists.
The zero-length list gives both units. The inclusions $d^R$ are
coassociative, and their zero-endpoint maps are identities.

For interchange, fix input charges $a,b$ and output charges
$p,q$, with $p+q=a+b$. The required right side is
\begin{equation}\label{gh:four-block}
 \sum_{\substack{i+j=p\\0\leq i\leq a,\ 0\leq j\leq b}}
 (m^R_{i,j}\otimes m^R_{a-i,b-j})
 (1\otimes\beta_{a-i,j}\otimes1)
 (d^R_{i,a-i}\otimes d^R_{j,b-j}).
\end{equation}
An $(a,b)$-interleaving has uniquely determined numbers $i,j$ of
its two input lists among the first $p$ output positions.
They are the first $i$ and first $j$ entries of those lists.
Exchange the middle blocks of lengths $a-i$ and $j$, then
interleave the two left pieces and the two right pieces.
This is a bijection with the terms in \eqref{gh:four-block}, and
the resulting permutations are identical. Their $R$-operators
are consequently identical. This proves interchange in every
charge, including all cases with a zero block.
Further refinements are the same partitions of an interleaving,
so their regrouping identities agree.

For $F\in\mathsf B_n$, full $R$-invariance gives
$m^R_{a,n-a}d^R_{a,n-a}(F)=\binom naF$.
Both antipode equations therefore reduce, for $n>0$, to
$\sum_{a=0}^n(-1)^a\binom na=0$; at zero they are the identity.
Here strict primitives mean the kernel, in positive charge, of
the reduced coproduct into pairs of positive charges.
Every charge-one element is primitive. For $n>1$ the component
$d^R_{1,n-1}$ is injective, so there are no further primitives.

In tensor degree $n$, the ideal in
\eqref{gh:braided-symmetric} is the span of
$(1-R_i)T_n$ for adjacent $i$.
It is the symmetric-group coinvariant relation. Over $\mathbb Q$,
averaging identifies coinvariants with invariants, and the stated
map is $n!e_n$, so it is an isomorphism onto $\mathsf B_n$.
Decomposing permutations into a shuffle and two internal block
permutations proves compatibility with concatenation and
\eqref{gh:abstract-operations}.
Equivalently, declare $N$ primitive in the braided tensor algebra.
For a quadratic relation the mixed coproduct is
$(1+R)(1-R)=0$. Thus the relation ideal is a Hopf ideal, and its
quotient has exactly the coproduct just constructed.
\end{proof}

\subsection{Comparison with the actual full flag trace}

The preceding construction was algebraic. The next theorem identifies
its maps with the geometric product and direct sum on the common
torsor.

\begin{theorem}[Global geometric Hopf comparison]\label{gh:global-hopf}
For the normalized rational smooth spherical coefficients
\eqref{gh:relative-coefficient}, the simultaneous comparison of
Theorem~\ref{gh:simultaneous} identifies the actual Hall product
\eqref{gh:hall-product} and actual direct-sum coproduct
\eqref{gh:direct-sum-map} with \eqref{gh:abstract-operations}.
It gives a coherent Hopf lift of these global common-$BE$ operations
to the tensor recipient $\mathcal T_R$. The crossing is the
polynomial relative operator \eqref{gh:polynomial-R}, and the
relative unit remains $\mathbb Q_{B_E}$.
\end{theorem}
\begin{proof}
Let $P_n:T_n\to\mathsf B_n$ be the actual iterated charge-one
Hall product in the simultaneous action model. By the binary
calculation,
$P_2=1+R$ after inclusion into $T_2$.
Associativity of the full flag product allows every adjacent pair
to be multiplied first, so $P_nR_i=P_n$ for each $i$.
Consequently $P_n=P_ne_n$.

Now take $F\in\mathsf B_n$. Under
\eqref{gh:curve-cohomology}, it is the direct-sum pullback of a
class on $\mathcal M_n$. On the full flag stack its
associated-graded pullback is therefore the pullback of that
same class along the proper map to $\mathcal M_n$.
The projection formula for the actual proper trace gives
\[
 P_n(F)=F\,P_n(1)=n!F,
\]
where the last equality is Proposition~\ref{gh:constants}.
This trace includes every nonsplit flag. Combining the two
equalities yields the operator identity
\begin{equation}\label{gh:total-symmetrizer}
 P_n=\sum_{w\in S_n}R_w
\end{equation}
on the entire $T_n$, in every charge.

For $F\in\mathsf B_a\otimes\mathsf B_b$, flag associativity gives
\[
 P_{a+b}(F)
 =m_{a,b}\bigl((P_a\otimes P_b)F\bigr)
 =a!b!\,m_{a,b}(F).
\]
Every permutation is uniquely a shuffle followed by an element of
$S_a\times S_b$. Since $F$ is fixed by that block subgroup,
\eqref{gh:total-symmetrizer} gives
$m_{a,b}(F)=\sum_{w\in\operatorname{Sh}(a,b)}R_wF$.
The actual direct-sum map was already identified with the inclusion
$d^R_{a,b}$ in Proposition~\ref{gh:direct-sum}.

Both identifications use the same natural monoidal formality
comparison, applied to the product, direct-sum, action and full flag
diagrams before descent. Thus they retain the actual coherences.
The interleaving identities of Theorem~\ref{gh:formal-hopf} supply
interchange and all its refinements in this common model, and
transport them back through that comparison.
\end{proof}

Theorem~\ref{sh:hopf} constructs these operations before
forgetting support, using the proper full-flag Springer action.
Its tensor recipient uses the universal permutation group
algebras. Corollary~\ref{sh:global-crossing} identifies the
global image of that action with the polynomial operator above.
This supplies a supported construction without requiring the
supported action to factor through the image quotients that
define $\mathcal T_R$.

If the same extension correspondence and coefficient map are recorded
quotient first, their product is $m^{\mathrm{quot}}=m\tau$.
No proper trace or coefficient is changed by this endpoint relabelling.
The crossing is correspondingly
\[
 R^{\mathrm{quot}}=\tau R\tau=\tau+\delta_{12}\partial_z.
\]
Reversing the full flag gives the coherent opposite-algebra comparison,
with the same direct-sum map. This is conjugation by the ordinary
endpoint swap, not replacement by $R^{-1}$, since $R^{-1}=R$.
A separate change of critical coefficient or orientation would require
its own comparison.

\begin{example}[The indispensable crossing]\label{gh:crossing-example}
Keep the Hall product fixed. Replacing only $R$ in the binary
intertwining equation by $\tau$ gives, on $z\otimes1$,
\[
 m\tau(z\otimes1)-m(z\otimes1)=2\delta_{12}\ne0.
\]
There is a four-block version. In input charges $(2,2)$ and output
charges $(2,2)$, use
$F=2(z_1+z_2-\delta_{12})$, representing
$(z*1)\otimes(1*1)$.
In only the $i=a-i=j=b-j=1$ term of \eqref{gh:four-block},
replace the middle $R_{23}$ by $\tau_{23}$.
The left side minus this altered right side is
\[
 -2(\delta_{13}+\delta_{14}+\delta_{23}+\delta_{24}).
\]
Indeed $(R_{23}-\tau_{23})F=-2\delta_{23}$.
The following two binary products act on this Chern-independent
class by $(1+\tau_{12})(1+\tau_{34})$, giving the displayed sum.
Setting $a_1=b_1=a_2=b_2=a_4=b_4=0$ leaves $-4a_3b_3$,
so the defect is nonzero. All other summands, including the whole
two-by-two block crossing, remain fixed in this comparison.
\end{example}

\subsection{The ordinary quotient after fibre realization}

Choose a trivial torsor $b:\mathrm{pt}\to BE$.
First form the actual derived fibre $b^*\mathcal A$, then take its
graded rational cohomology. The result is $\mathcal H=\mathsf B$
with the operations just compared. Its cohomological degrees are
nonnegative.

\begin{corollary}\label{gh:fibre-quotient}
The subspace
\[
 J=\bigoplus_{n\geq0,\ j>0}H^j(\mathcal M_n;\mathbb Q)
\]
is an ideal and a coideal for the specified Hopf operations, and is
stable under the crossings and antipode. The quotient is the
ordinary even polynomial Hopf algebra $\mathbb Q[x]$, with $x$
primitive of charge one. Its projection and section are
\[
 \rho(\mathbf1_n)=\frac{x^n}{n!},\qquad
 \rho(J)=0,\qquad
 \jmath(x^n)=n!\mathbf1_n,\qquad \rho\jmath=1.
\]
The actual direct-sum coproduct satisfies
\[
 d(\mathbf1_n)=\sum_{a+b=n}\mathbf1_a\otimes\mathbf1_b.
\]
The kernel is nonzero in every positive charge.
\end{corollary}
\begin{proof}
All operations preserve total cohomological degree. Positive degree
is an algebra ideal, and every positive total degree in a tensor
product has a positive-degree factor, proving the coideal assertion.
Degree-zero projection is monoidal on nonnegatively graded vector
spaces. On degree zero, the divided difference in $R$ vanishes,
so this functor carries the specified crossings to the ordinary
even swap. It therefore carries the Hopf diagrams to ordinary
Hopf diagrams.

The constant-class multiplication in
Proposition~\ref{gh:constants} gives the normalization of $\rho$
and $\jmath$. Direct-sum pullback of a constant is constant, giving
the coproduct formula. Under $\rho\otimes\rho$, its right side is
\[
 \sum_{a+b=n}\frac{x^a}{a!}\otimes\frac{x^b}{b!}
 =\Delta\left(\frac{x^n}{n!}\right).
\]
The section is likewise a Hopf map.
For every $n>0$, the nonzero invariant
$z_1+\cdots+z_n$ belongs to $J\cap\mathsf B_n$.
\end{proof}

The quotient is taken after fibre realization and graded cohomology.
Before forgetting the action, $J$ is not a $\Lambda_E$-submodule:
$I_a(a)=1$. Thus Corollary~\ref{gh:fibre-quotient} is compatible
with Proposition~\ref{gh:no-relative-rho}; it supplies no
normalized map to $\mathbb Q_{B_E}[x]$.
It recognizes an ordinary one-generator positive enveloping algebra
inside the fibre realization. It does not identify the full Hall
carrier with that algebra or recover a Cartan matrix from it.

\begin{corollary}[Charge windows and completion]
\label{gh:windows}
All the displayed operations and the fibre quotient restrict
compatibly to total-charge windows $0\leq n\leq N$.
Their tensor at charge $n$ is
$\bigoplus_{a+b=n}V_a\otimes W_b$, with $n\leq N$.
For rectangular output bounds $r,s$, the required input bound is
$r+s$.
After the ordinary fibre realization, the charge completion
\[
 \widehat{\mathcal H}=\prod_{n\geq0}\mathsf B_n,\qquad
 \widehat{\mathcal H}\,\widehat\otimes\,\widehat{\mathcal H}
       =\prod_{a,b\geq0}(\mathsf B_a\otimes\mathsf B_b)
\]
has continuous product, coproduct, block crossings and antipode.
It has a split continuous quotient
$\widehat\rho:\widehat{\mathcal H}\to\mathbb Q[[x]]$ with kernel
$\prod_n(J\cap\mathsf B_n)$.
\end{corollary}
\begin{proof}
Every component has its stated total charge, so deleting components
of charge greater than $N$ commutes with all remaining component
maps. This gives the compatible restriction maps and sections.
An output pair $(a,b)$ uses only input charge $a+b$, proving the
rectangular bound.
Give each $\mathsf B_n$ the discrete topology and the products the
product topology. A fixed multiplication component uses the finitely
many pairs $a+b=n$; a fixed coproduct component uses only
$\mathsf B_{a+b}$. The same property holds for crossings, antipode
and quotient. They are therefore continuous, and all their
identities follow componentwise from the finite-charge identities.
The kernel and section are likewise componentwise.
\end{proof}

This is only completion in charge. In any one charge, a cohomology
class still contains finitely many Chern monomials. Neither a
Chern-degree completion nor a support or distributional completion
has been introduced. The total-charge tensor also differs from the
ordinary tensor square of the quotient discarding charges greater
than $N$; the latter retains output pairs whose sum exceeds $N$.
